% =====================================================================
% Results — RESTRUCTURED DRAFT for the full confirmation sweep (E2).
%
% NOT part of the build. `paper/sec_results.tex` remains the submittable
% version and stays authoritative until E2 lands and this file replaces it.
%
% \PH{...} marks a value awaiting E2. The wording guard skips *_draft.tex but
% flags \PH{ anywhere else, so promoting this file with an unfilled placeholder
% fails the check rather than shipping.
%
% Spine: full confirmation ranking -> how far the screening protocol displaces
% it -> three sources of sparsity -> delta reduces to thresholding. The second
% subsection is the paper's central claim, measured rather than argued.
%
% Populations, which must never blur:
%   screening        27 configs x 20 blocks  (4 datasets x 5 seeds, v1)
%   confirmation     27 configs x 25 blocks  (5 protocols x 5 seeds, v2)
%   encoding axis     3 encodings x 25 blocks   (registered)
%   neuron axis       9 neurons x 20 blocks     (registered)
% =====================================================================

\section{Results and Discussion}

\subsection{The design space under the confirmation protocol}
\label{sec:res-conf}

% results/analysis/v2_full/{mean_rank.csv,friedman.json}
Under the leakage-resistant confirmation protocol, all \FNScrVariants{}
configurations were evaluated on \PH{5} protocols at five seeds, giving
\PH{25} paired (protocol $\times$ seed) blocks. \PH{Leader} attains the lowest
mean macro-F1 rank (\PH{rank}); \texttt{LeakyParallel/latency} places
\PH{position} at \PH{rank}. Friedman rejects equal performance
($\chi^2=\PH{chi}$, $p=\PH{p}$).

The Nemenyi critical difference at $k=\FNScrVariants$ and $N=\PH{25}$ is
\PH{cd}. As registered in advance, we report the post-hoc descriptively: with
this many configurations and this few blocks the all-pairs resolution is coarse,
so a failure to separate two configurations is a statement about the test's
power and not evidence that they perform equally. The comparisons that carry
weight are the pre-registered paired ones and the protocol-level collapse below.

% results/analysis/protocol_level/
Collapsing each protocol to a single mean before ranking, which treats the
protocol rather than the seed as the unit and so does not count five seeds on
one dataset as five independent observations, gives \PH{protocol-level
summary}.

\subsection{How far the screening protocol displaces the ranking}
\label{sec:res-shift}

% results/analysis/protocol_shift/tau.json -> fig_rank_shift
This is the question the study exists to answer, and it reduces to one number.
The same \FNScrVariants{} configurations ranked under the two protocols give a
Kendall rank correlation of $\tau=\PH{tau}$ (permutation $p=\PH{taup}$,
\PH{10000} permutations at the seed fixed in the registration addendum).
Figure~\ref{fig:rankshift} shows the movement configuration by configuration.

\PH{Sentence naming the configurations that move most, in each direction.}

% results/v1_submitted/ctu13_strict_vanilla.csv
The displacement is not uniform across datasets, and CTU-13 shows where it comes
from. Under the screening protocol, with whole-scenario host aggregates and a
random flow split, \texttt{LeakyParallel/latency} attains \FNCtuLegacy{}
macro-F1. Making the aggregation causal alone leaves \FNCtuCausalRandom.
Adding a scenario-disjoint partition brings it to $\mathbf{\FNCtuStrict\%}$,
with a standard deviation of \FNCtuSdEval{} percentage points over the
\FNCtuNEval{} fold-seed evaluations and fold means, in fold order, of
\FNCtuIdxFA, \FNCtuIdxFB, \FNCtuIdxFC{} and \FNCtuIdxFD. A single benchmark
therefore accounts for much of the ranking movement, and it does so because the
screening partition let scenario- and host-related information cross the
evaluation boundary.

\subsection{Registered confirmation along the two design axes}
\label{sec:res-axes}

The analysis registered before any confirmation data existed varies one factor
at a time. It remains the confirmatory result; the full sweep of
Section~\ref{sec:res-conf} is more complete but was registered only ahead of its
own runs, after these results were known, and we report both.

Encoding axis, neuron held at \texttt{LeakyParallel}, \FNEncBlocks{} blocks:
latency attains mean rank \FNEncRank{} of \FNEncLevels{} and wins \FNEncWins{}
blocks (Friedman $p=\FNEncP$). Neuron axis, encoding held at latency,
\FNNeuBlocks{} blocks: \texttt{LeakyParallel} attains \FNNeuRank{} of
\FNNeuLevels{} ($p=\FNNeuP$), narrowly, with \FNNeuRunnerUp{} at
\FNNeuRunnerUpRank{} and \FNNeuMostWins{} winning more individual blocks.

\PH{If the full sweep and the registered axes disagree, state it here per
addendum 6.1: report both, name both leaders, do not adjudicate.}

\subsection{Three sources of sparsity}
\label{sec:res-sparsity}

% carried forward unchanged from the current sec_results.tex
\PH{Import the existing subsection: zero-feature mass, the input-ratio identity
verified to \FNPredErrMax\%, the hidden-layer dilution, and the SOP result.}

\subsection{Delta reduces to one-step thresholding}
\label{sec:res-delta}

% carried forward unchanged
\PH{Import the existing derivation and the temporal-convergence table.}

\subsection{Robustness and boundary conditions}
\label{sec:res-limits}

% carried forward unchanged
\PH{Import deduplication, CIC-IDS2017 U2R support, threshold transfer, and
known-versus-unseen.}
